\section{A companion equation in ten unrestricted integers}
\label{sec:sun-companion}

The same quadratic coefficient packing also compresses the ten-integer
construction already present in source~31 of the repository~\cite{repo}.
The resulting polynomial has
\[
 \boxed{10\text{ integer unknowns},\qquad
 \deg F_{10}=449\,607\,016,\qquad 773\text{ binary operations}.}
\]
The earlier source~31 equation has degree
$589\,042\,901\,992\,797\,851\,562\,527\,064$ and 1,442 operations.
The comparison uses the same convention: every binary addition,
subtraction and multiplication is charged, including multiplication by
a fixed coefficient; repeated registers and fixed leaves are free.
Thus this application improves both degree and circuit size, while
retaining the same ten unrestricted integer unknowns.  It adds the
Lucas-theoretic dependency on Sun's 2026 preprint~\cite{sun2026}; the
nine-natural-variable construction does not need that preprint.

\subsection{A different source scale, with the same quadratic kernel}

Use the same 31 shifted quadratic coordinates and 22 residuals $r_j$,
the same Sidon positions $0,a_1,\ldots,a_{31}$, and the same
$M=799$, $d=3197$, $H=70334$ as in the preceding packing construction.
With $|\alpha|\le2$, write
\[
 r_j(\mathbf n)=\sum_\alpha p_{j,\alpha}\mathbf n^\alpha,
 \qquad c_{j,\alpha}=(2-|\alpha|)!\prod_i\alpha_i!\,p_{j,\alpha}.
\]
Define the fixed polynomial and target mask
\begin{align*}
 E(T)&=\sum_{j=0}^{21}\sum_{|\alpha|\le2}
 c_{j,\alpha}T^{jd+2M-\sum_i\alpha_i a_i},\\
 Z(T)&=\sum_{j=0}^{21}T^{jd+2M}.
\end{align*}
There are 136 nonzero weighted coefficient records and 16 distinct
nonliteral coefficient recipes. All four native ports are bound to the
new tuple compiled from complemented labels in the sibling \texttt{compiler} directory.

For each monomial $q\prod_i\mathrm{port}_i^{e_i}$ in these records,
majorize its magnitude by $|q|2^{\sum_i e_i b_i}$, where $b_i$ is the
actual bit length of the corresponding native port. Let $C_{\mathrm S}$ be one
plus the sum of all these majorants. This bit-length majorant is distinct
from the exact positive-monomial majorant used for the nine-variable source. The delivered exact arithmetic gives
\[
 \operatorname{bitlength}(C_{\mathrm S})=1\,359\,515.
\]
Set
\[
 \beta_{\mathrm S}=2^{1\,359\,528}>2048C_{\mathrm S}.
\]
Its exponent is even, so $\beta_{\mathrm S}\equiv1\pmod3$.
The binary payload of this additional constant occupies 169,942 bytes
before compression and 200 bytes in the delivered compressed form.
The number is fixed by the finite compiler data, independently of
whether a refutation exists.

For integer variables $f,g$, define
\begin{align*}
 b&=1+3f,& \mathcal B&=\beta_{\mathrm S}b^2,&
 \ell&=\sum_{i=1}^{31}\mathcal B^{a_i},\\
 \lambda&=\sum_{t=0}^{H-1}\mathcal B^t,&
 N_0&=\mathcal B^{M+1},& N_1&=4\mathcal B^H,\\
 N&=N_0N_1=4\mathcal B^{71134},&
 \mathscr C&=b\mathcal B^M,&
 M_0&=N_0-1-(b-1)\ell.
\end{align*}
Put
\[
 J_0=2(1+g)^2E(\mathcal B)+\mathcal B\lambda,
\]
and then
\begin{align*}
 S_0&=g+N_0J_0,&
 T_0&=M_0+N_0(\mathcal B-2)Z(\mathcal B),\\
 R_0&=(S_0+T_0+1)N+T_0+1,&
 X&=(N-1)R_0,&Y&=N^2.
\end{align*}
Here $E$ is a polynomial abbreviation rather than a new unknown.
All powers and summation limits are fixed. No sign restriction is
imposed on these polynomials for arbitrary free inputs: the reverse
bridge first recovers the bounds used in the next paragraph.

\paragraph{The source contract.}
These formulas have the following three properties.
If the quadratic system has a natural solution, there are arbitrarily
large $f>0$ and a $g$ with $b\le g<\mathscr C$, such that $b$ is a
square and a power of two, and $Y\mid\binom{2X}{X}$.
Conversely, $f\ne0$, $0\le g<2\mathscr C$, the same two conditions on
$b$, and that divisibility imply a natural solution of the system.
Finally, whenever $f\ne0$, $b$ is a square and $0\le g<2\mathscr C$,
one has $X>3b$ and $Y>\max\{b,256\}$, even before imposing the
power-of-two condition.

\paragraph{Proof of the changed source step.}
A square $b=1+3f$ with $f\ne0$ has $b\ge4$ and $f>0$.
The highest exponent of $E$ is at most $H-2M-1$, so
\[
 |E(\mathcal B)|<2C_{\mathrm S}\mathcal B^{H-2M-1}.
\]
For $0\le g<2b\mathcal B^M$, the inequality
$1+g\le3b\mathcal B^M$ gives
\[
 |2(1+g)^2E(\mathcal B)|
 <36C_{\mathrm S}b^2\mathcal B^{H-1}<\mathcal B^H.
\]
Since $\mathcal B^H\le\mathcal B\lambda<2\mathcal B^H$,
we obtain $0<J_0<3\mathcal B^H<N_1$ and $g<N_0$.
The mask $M_0$ has digits $\mathcal B-b$ at the positions $a_i$
and $\mathcal B-1$ elsewhere, so $0\le M_0<N_0$.
Also $(\mathcal B-2)Z(\mathcal B)<N_1$.
Consequently $0\le S_0,T_0<N$, $R_0\ge N+1$,
$X\ge N^2-1>3b$, and $Y>\max\{b,256\}$.
Moreover $N\equiv1\pmod3$, so $3\mid X$.

If $b$ is a power of two, so are $\mathcal B,N_0,N_1,N$.
Sun's central-binomial carry lemma~\cite[Lemma 2.2]{sun2021} gives
\[
 Y\mid\binom{2X}{X}
 \quad\Longleftrightarrow\quad
 \tau_2(S_0,T_0)=0.
\]
The two blocks split this into the code-validity test
$\tau_2(g,M_0)=0$ and the residual test
$\tau_2(J_0,(\mathcal B-2)Z(\mathcal B))=0$.
The first test decodes
$g=\sum_i n_i\mathcal B^{a_i}$ with $0\le n_i<b$ by
\cite[Lemma 2.4]{sun2021}, which explicitly permits arbitrary distinct
digit positions. For such digits, put
\[
 K_0=(1+g)^2E(\mathcal B)+(\mathcal B/2)\lambda,
 \qquad J_0=2K_0.
\]
Every coefficient of $K_0$ lies strictly between 0 and $\mathcal B$,
because its signed correction has magnitude at most
\[
 C_{\mathrm S}\left(1+\sum_i n_i\right)^2
 <1024C_{\mathrm S}b^2<\mathcal B/2.
\]
The selected digits are $\mathcal B/2+2r_j(\mathbf n)$.
Adding $(\mathcal B/2-1)Z(\mathcal B)$ without a binary carry
forces each selected digit to equal $\mathcal B/2$, hence every
$r_j(\mathbf n)=0$. This proves soundness.
Conversely, a simultaneous zero is nonzero because the fifth residual
at $\mathbf n=0$ equals 1. Choose $b=4^q$ larger than all its coordinates,
set $f=(b-1)/3$, and encode it in $g$. Then $b\le g<\mathscr C$,
and both carry tests hold. This proves completeness, with $q$ arbitrarily
large. No sum of squares is used in this source step.

\subsection{The precise imported Lucas interface}

The positive and reverse bridge proofs in
\cite[Propositions 6.1 and 6.2]{sun2026} use precisely the three-part
source contract of its Proposition 2.1, established above for the new
source, the fixed form $b=1+3f$, and the monomial bound
$\mathscr C=2^{(1\,359\,528)M}b^{2M+1}$.
Their proofs do not otherwise use the old base-five coefficient code.
We use this proof-level specialization, rather than asserting that the
new source is literally Sun's original source polynomial.
The positive bridge supplies a Lucas certificate, one square and one
divisibility condition; the reverse bridge recovers the source contract
from those conditions and a positivity inequality.
Sun's modified Lucas certificate~\cite[Theorem 3.1]{sun2026}
and full residue interpolation~\cite[Theorem 4.1]{sun2026} are additional
imported results. The latter permits the second Lucas witness to lie
in any prescribed residue class modulo a suitable first witness.
This permits a positivity witness and the divisibility
merger to share variables. The positive bridge is the specialization
$p=P=2$, $Q=1$ of \cite[Theorem 4.12]{sun2021}; the reverse bridge
uses the estimates of \cite[Theorem 4.13]{sun2021}, replacing its
Lucas-certificate call by \cite[Theorem 3.1]{sun2026}.

The complete ten-variable list is
\[
 (f,g,h,k,l,w,x,y,z,m)\in\Z^{10}.
\]
Use the source quantities above and define
\begin{align*}
 L&=lY,&U&=2LX,\\
 V&=4gwY,&W&=bw,\\
 A&=U(V+1),&B_1&=2X+1,\\
 C_1&=B_1+(A-2)h,&P&=U^2V,\\
 K&=X+1+k(P-2),&S&=2A-5,\\
 T&=3WC_1-2(W^2-1),&r_1&=(P^2-4)K^2+4.
\end{align*}
Next set
\begin{align*}
 \Delta&=8\mathscr C^3gK^2
       -g^2\bigl(32(C_1-KL)^2\mathscr C^3+g^2K^2\bigr),\\
 O&=f^2l^2\Delta,\qquad\xi=zx.
\end{align*}
The modified Lucas quantities are
\begin{align*}
 D&=(A^2-4)C_1^2+4,& E_1&=C_1^2D\xi,\\
 F&=4(A^2-4)E_1^2+1,&G&=2(A+2)(A-2)^2E_1^2-1,\\
 \Omega_1&=16E_1,&H_1&=C_1+(B_1-C_1)F+\Omega_1Fy,\\
 I&=(G^2-1)H_1^2+1,&r_2&=DFI,\\
 r_0&=b\bigl(1+(3O-4)b^2z^2\bigr).&&
\end{align*}
The positive bridge has $O\ge4$, since $g\ge b\ge4$ and its positive
integer $\Delta$ is divisible by $g$.
As proved in source~31~\cite{repo}, the coprimality of $b$ and
$1+(3O-4)b^2z^2$, together with $b\equiv1\pmod3$, makes $r_0$ square
and $z\ne0$ enforce $b$ square and $O>0$. Conversely, a Pell solution
with its second coordinate divisible by $b$ supplies such $z$.

\subsection{The final polynomial and its implication in both directions}

Let
\begin{align*}
 \pi(v)&=(v^3-v-1)(v^2+23),&\Sigma&=\xi S^5,\\
 t&=T+Sy,&
 \Theta&=(t^3-tS^2-S^3)(t^2+23S^2).
\end{align*}
Define the weights and scaled radicands by
\[
 w_0=1,\qquad w_1=r_0,\qquad w_2=r_0(P^2+K^2+3),
\]
\[
 Z_0=\Sigma m-\Theta,\qquad
 A_i=\Sigma^2w_i^2r_i\quad(0\le i\le2),
\]
and put
\[
 V_0=Z_0^2+A_2-A_0-A_1,\qquad
 W_0=V_0^2+4Z_0^2A_2-4A_0A_1.
\]
The requested integer equation is
\[
 \boxed{F_{10}=W_0^2-16V_0^2Z_0^2A_2=0.}
\]
Every symbol other than the ten listed inputs is a polynomial
abbreviation or a fixed integer.

The local quintic merger and norm arguments are inherited from
source~31~\cite{repo}. Byte-exact copies of that proof, its mathematical
review and its finalizer are supplied under \texttt{signed10/inherited/};
the historical coefficient binding there is inactive. For clarity, their exact roles are as follows.
The polynomial $\pi$ has no rational root and has a root modulo every
positive integer. Thus $\Theta\ne0$, since $S=2A-5\ne0$.
Divisibility $\Sigma\mid\Theta$ forces $\xi\ne0$ and $S\mid T$:
if $v_p(S)>v_p(T)$, the two homogeneous factors have valuations
$3v_p(T)$ and $2v_p(T)$, contradicting divisibility by $S^5$.
Conversely, if $S\mid T$ and $\xi\ne0$, a suitable residue class of $y$
modulo $|\xi|$ makes the divisibility true. Full residue interpolation
selects that class while preserving the Lucas square certificate.
This converse is existential in $y$, not a pointwise assertion for every $y$.

The norm identity is
\[
 F_{10}=\prod_{\varepsilon\in\{\pm1\}^3}
 \left(Z_0-\Sigma\sum_{i=0}^2\varepsilon_iw_i\sqrt{r_i}\right).
\]
At a zero, $\Sigma=0$ would imply $F_{10}=\Theta^8\ne0$, so
$z\ne0$. Then $|r_0|>1$, and the structural weights satisfy
\[
 |w_1|>\sqrt{|r_0|},\qquad
 |w_2|>\sqrt{|r_0|}+|w_1|\sqrt{|r_1|}.
\]
A rational signed sum of these roots must have each root rational:
a Galois automorphism changing the largest indexed root cannot cancel
its change with the smaller terms. Rational square roots of integers
are integers, so the three radicands are squares and $\Sigma\mid\Theta$.
This reasoning also covers negative candidate radicands, zero roots,
and repeated square classes.

The reverse Lucas bridge therefore gives a simultaneous natural zero
of the 22 residuals. Conversely, its positive bridge, the coprime Pell
step and full residue interpolation give the three square conditions
and $\Sigma\mid\Theta$; an integer $m$ then kills one norm factor.
Together with the native compiler and calibrated K215 endpoint, this
proves
\[
 \mathrm{ZFC}\vdash
 \bigl(\neg\exists\mathbf x\in\Z^{10}\ F_{10}(\mathbf x)=0\bigr)
 \longleftrightarrow\operatorname{Con}(\mathrm{ZFC}).
\]
This is an equivalence based on the stated mathematical dependencies,
not a proof of consistency or a new proof-assistant certificate.

\subsection{Degree and actual-instance checks}

The degrees of $X,Y$ are 426,804 and 284,536. The three scaled radicands
$A_0,A_1,A_2$ have degrees 13,947,118, 24,764,324 and 112,401,754.
The last dominates the squared norm center and the other radicands, so
the final leading part is $A_2^4$. The claimed exact degree is therefore
$4\cdot112\,401\,754=449\,607\,016$.
Independently, diagonal leading coefficients for the actual complemented
tuple are nonzero modulo 101, 1009 and 1,000,003: they are 81, 287 and 270,876.
These are exact finite certificates that the structural degree bound is
attained after specializing all fixed coefficients.

The verifier independently reconstructs the 136 records by evaluating
quadratic values, gradients and homogeneous parts at the all-one
coordinate vector, while the builder uses binomial expansion. It checks
all 528 pair sums, every native-port hash and strict coefficient bound,
the materialized $\beta_{\mathrm S}$, the formal norm expansion, and
1,620 named-register comparisons against literal scalar formulas over
three finite fields. These checks establish the finite artifact's
transcription and quantitative claims; the unbounded equivalence rests
on the proofs above and the stated imported theorems.
